\section{结论、优缺点与推广}

问题一的确定性混合整数线性规划给出了满足功率平衡、充放电互斥和储能边界的日内调度方案。
问题二进一步用 50 个成对负荷--光伏场景描述预测误差，并严格区分日前模型评价与事后实际回测。
334 个日模型均取得最优解，模型内期望总成本为 $14\,374\,345.246186$ 元；按附件 1 固定日内
电价制定计划、以附件 2 实际负荷和实际光伏回测后，实际紧急购电量为
$398\,161.900954\ \mathrm{kWh}$，实际总购电费用为 $14\,839\,462.002357$ 元。

该模型的优势是场景配对保持了负荷与光伏误差的相关性，实际执行层又以计划充放电为上限，
通过实际动作递推 SOC 并跨日传递，从而避免未吸收放电等无物理解释的能量渠道。其局限主要是
场景数量与历史残差样本有限，预测模型也尚未利用更丰富的天气信息。

问题三在相同物理契约上引入多时刻预测与滚动优化。334 天回测中，加入 6:00 与 12:00
信息后实际总费用由 $1472.797$ 万元降至 $1450.550$ 万元，相对 $S_0$ 的信息价值为
$22.247$ 万元；继续加入 18:00 信息后总费用为 $1450.772$ 万元，增量信息价值为
$-0.222$ 万元。因此，本次事后回测中 $S_2=\{0,6,12\}$ 的经济表现最好，但不据此宣称其为
事前最优安排；$S_3$ 完整策略保留在附件中供核验。多次调整采用相对上一承诺的主结算口径，
同时以相对 00:00 承诺口径验证结论稳定性。
